\documentclass[12pt]{article}
\usepackage[utf8]{inputenc}
\usepackage[margin=0.75in]{geometry}
\usepackage{amsmath}\usepackage{amssymb}\usepackage{enumitem}\usepackage{multicol}\usepackage{tcolorbox}
\setlength{\parindent}{0pt}
\begin{document}

\begin{center}{\Large \textbf{Identify parent functions and their key features}}\\[2pt]{\small Pre-Cal \textbar\ CK-12: Function Families}\end{center}\vspace{0.25cm}

{\large\textbf{Key Idea}}\par\smallskip
A \textbf{parent function} is the simplest member of a family — no shifts, stretches, or reflections. Learn each shape and one key point.\par


{\large\textbf{Key Terms}}\par\smallskip
\begin{itemize}[leftmargin=1.4em]
  \item \textbf{Parent function} — the base graph of a family
  \item \textbf{Domain} — allowed $x$; \textbf{Range} — resulting $y$
\end{itemize}


{\large\textbf{Worked steps}}\par\smallskip
Step 1: Read the rule and name the family.\par
Step 2: Recall its shape, key point, and domain.\par
Example: $f(x)=\sqrt{x}$ is the \textbf{square-root} parent; it starts at $(0,0)$ and only exists for $x\ge 0$.\par


{\large\textbf{You Try}}\par\smallskip
Name the parent of $f(x)=\dfrac{1}{x}$ and state its domain.\par

\end{document}